\chapter{Правова база протидії кіберзлочинності в Україні}
\label{apdx:legal}

\noindent\textit{Довідкова інформація для швидкого реагування та оцінки правових наслідків.}

\vspace{0.5cm}

Кримінальний кодекс України встановлює конкретну відповідальність за кіберзлочини. У таблиці Б.1 наведено основні статті, які застосовуються при розслідуванні випадків онлайн-шахрайства, несанкціонованого доступу до комп'ютерних систем та розповсюдження шкідливого програмного забезпечення.

Важливо розуміти, що злочинець не завжди уникає покарання через час. Строки давності залежать від тяжкості злочину, а також існують офіційні канали для термінового звернення. У таблиці Б.2 наведено інформацію про строки давності притягнення до кримінальної відповідальності та контактні дані державних органів, куди можна звернутися у разі виявлення кіберзлочину.



\section*{Таблиця Б1. Відповідальність за кіберзлочини (ККУ)}

\begin{table}[H]
    \centering
    \renewcommand{\arraystretch}{1.4}
    \begin{tabularx}{\textwidth}{@{} l X X @{}}
        \toprule
        \textbf{Стаття} & \textbf{Суть порушення} & \textbf{Максимальне покарання} \\
        \midrule
        
        \textbf{Ст. 190} & 
        \textbf{Шахрайство.} Заволодіння майном/правом на нього шляхом обману чи зловживання довірою. & 
        Штраф або \textbf{позбавлення волі до 12 років} (з конфіскацією). \\
        \cmidrule{1-3}
        
        \textbf{Ст. 361} & 
        \textbf{Несанкціоноване втручання.} Злам систем, що призвів до витоку, втрати чи блокування даних. & 
        Обмеження або \textbf{позбавлення волі від 3 до 6 років}. \\
        \cmidrule{1-3}
        
        \textbf{Ст. 361-1} & 
        \textbf{Шкідливе ПЗ.} Створення та розповсюдження вірусів, троянів, експлойтів. & 
        Виправні роботи або \textbf{позбавлення волі до 3 років}. \\
        \cmidrule{1-3}  

        \textbf{Ст. 361-2} & 
        \textbf{Збут кодів доступу.} Продаж/передача паролів для зламу систем. & 
        Штраф або \textbf{позбавлення волі до 2 років}. \\
        \cmidrule{1-3}
        
        \textbf{Ст. 362} & 
        \textbf{Крадіжка даних.} Викрадення або привласнення комп'ютерної інформації. & 
        Виправні роботи або \textbf{позбавлення волі до 3 років}. \\
        \cmidrule{1-3}
        
        \textbf{Ст. 163} & 
        \textbf{Порушення таємниці.} Читання чужого листування, прослуховування дзвінків. & 
        Для службових осіб — \textbf{позбавлення волі від 3 до 7 років}. \\
        
        \bottomrule
    \end{tabularx}
\end{table}

\vspace{0.5cm}

\section*{Таблиця Б2. Строки давності та корисні контакти}

\begin{table}[H]
    \centering
    \renewcommand{\arraystretch}{1.3}
    \begin{tabularx}{\textwidth}{@{} X p{0.45\textwidth} @{}}
        \toprule
        \multicolumn{2}{c}{\textbf{Строки давності (Ст. 49 ККУ)}} \\
        \textit{Час, після якого особа звільняється від відповідальності} & \textit{Тип злочину} \\
        \midrule
        \textbf{2 роки} & Невеликої тяжкості \\
        \textbf{3 роки} & Середньої тяжкості (до 5 років тюрми) \\
        \textbf{5 років} & Тяжкі злочини (до 10 років тюрми) \\
        \textbf{10 років} & Особливо тяжкі (понад 10 років) \\
        \bottomrule
        \addlinespace[0.5cm] % Відступ між частинами таблиці
        
        \toprule
        \multicolumn{2}{c}{\textbf{Куди звертатися (Гарячі лінії)}} \\
        \midrule
        \textbf{Кіберполіція України} & \texttt{cyberpolice.gov.ua} \\
        \small{(Шахрайство, віруси, злами)} & Тел: \textbf{0-800-505-170} або \textbf{102} \\
        \midrule
        \textbf{Національний банк (НБУ)} & \texttt{bank.gov.ua} \\
        \small{(Фінансові махінації)} & Тел: \textbf{0-800-505-240} \\
        \midrule
        \textbf{Служба безпеки (СБУ)} & \texttt{ssu.gov.ua} \\
        \small{(Держзрада, кібертероризм)} & Тел: \textbf{0-800-501-482} \\
        \bottomrule
    \end{tabularx}
\end{table}

\noindent\footnotesize{\textit{* Увага: Законодавство змінюється. Актуальні норми на сайті \texttt{zakon.rada.gov.ua}}}